\documentclass[preprint,12pt,authoryear]{elsarticle}

\usepackage{amsmath,amssymb,amsfonts}
\usepackage{algorithm}
\usepackage{algorithmic}
\usepackage{booktabs}
\usepackage{multirow}
\usepackage{graphicx}
\usepackage{xcolor}
\usepackage{hyperref}
\usepackage{url}
\usepackage{subcaption}

% Custom commands
\newcommand{\sdg}{\textsc{SDG-MetaBBO}}
\newcommand{\R}{\mathbb{R}}
\newcommand{\E}{\mathbb{E}}
\newcommand{\xbest}{x^*}
\newcommand{\sigmaens}{\sigma_{\text{ens}}}

\journal{Applied Soft Computing}

\begin{document}

\begin{frontmatter}

\title{SDG-MetaBBO: Surrogate-Disagreement-Guided Differential Evolution with Goal-Distancing Operators for Black-Box Optimization}

\author[yalova]{Murat G\"ok\corref{cor}}
\ead{murat.gok@yalova.edu.tr}
\cortext[cor]{Corresponding author.}

\affiliation[yalova]{organization={Department of Computer Engineering, Faculty of Engineering},
            addressline={Yalova University},
            city={Yalova},
            postcode={77200},
            country={Turkey}}

\begin{abstract}
This paper presents \sdg{}, a meta-black-box optimization framework that integrates three mechanisms into a SHADE-based differential evolution backbone: (1)~a sliding-window UCB bandit that selects from six operators using rank-based credit augmented by surrogate ensemble disagreement, (2)~a goal-distancing operator that generates trial vectors moving \emph{away} from the current best with uncertainty-scaled step size and L\'evy flight for escaping deceptive basins, and (3)~a PPO-trained meta-controller that tunes exploration--exploitation parameters online. The surrogate ensemble (five independently initialized MLPs) serves as a learned landscape model whose prediction disagreement identifies uncertain regions---acting as both an exploration reward signal for the bandit and a step-size modulator for the goal-distancing operator. On six benchmark functions, \sdg{} with fixed meta-parameters achieves statistically significant improvements over L-SHADE on five of six functions (Wilcoxon $p < 0.001$, Vargha--Delaney $A_{12} > 0.77$, 30~independent runs), with median error reductions ranging from $1.1\times$ (Lunacek bi-Rastrigin) to $10{,}000\times$ (Rastrigin). The design is biologically motivated by the cognitive hunting strategies of the \emph{Portia} spider but defended entirely in optimization-theoretic terms.
\end{abstract}

\begin{keyword}
meta-black-box optimization \sep differential evolution \sep surrogate-assisted optimization \sep adaptive operator selection \sep goal-distancing \sep ensemble disagreement
\end{keyword}

\end{frontmatter}

%% ============================================================
\section{Introduction}
\label{sec:intro}
%% ============================================================

Black-box optimization (BBO)---minimizing an unknown objective $f(\mathbf{x})$ given only pointwise evaluations---is a core computational challenge in engineering design, hyperparameter tuning, and scientific discovery. Differential evolution (DE) \citep{Storn1997} and its adaptive variants L-SHADE \citep{Tanabe2014} and jSO \citep{Brest2017} remain among the most competitive solvers for continuous BBO, consistently performing well in CEC competitions. However, their fixed operator pipelines and heuristic parameter schedules cannot adapt to the heterogeneous landscape structures encountered across diverse problem instances.

The Meta-Black-Box Optimization (MetaBBO) paradigm \citep{Ma2025survey} addresses this limitation by formulating algorithm design as a bi-level problem: a meta-level controller $\pi_\omega$ selects operators, parameters, or entire algorithmic structures for a low-level BBO solver, with the meta-objective being expected performance across a problem distribution. Recent MetaBBO methods---RL-DAS \citep{Guo2024}, RLDE-AFL \citep{Guo2025}, ConfigX \citep{Guo2024configx}---have demonstrated that reinforcement learning (RL) controllers can outperform hand-tuned algorithms on standard benchmarks.

Despite this progress, three structural limitations persist in the current MetaBBO landscape:

\begin{enumerate}
\item \textbf{Operator credit ignores model uncertainty.} Existing RL-based operator selectors reward operators based on raw fitness improvement, which conflates ``genuinely effective operator'' with ``lucky evaluation in unexplored space.'' Incorporating surrogate model disagreement into the credit signal provides the missing distinction.

\item \textbf{No first-class operator for escaping deceptive basins.} All mainstream DE operators generate trial vectors that move \emph{toward} promising regions. On deceptive landscapes with multiple basins (e.g., Lunacek bi-Rastrigin), this directionality traps the optimizer in suboptimal funnels.

\item \textbf{Static meta-parameters across optimization phases.} The optimal exploration coefficient, goal-distancing intensity, and surrogate update frequency change as optimization progresses, but existing methods use fixed schedules.
\end{enumerate}

This paper proposes \sdg{} (Surrogate-Disagreement-Guided Meta-Black-Box Optimizer), which addresses all three limitations through a principled integration of adaptive operator selection, learned uncertainty, and deliberate anti-convergence moves.

\subsection{Biological motivation}

The design of \sdg{} is motivated by the cognitive hunting strategies of the \emph{Portia} spider (Salticidae), documented by \citet{Jackson2011}. \emph{Portia} exhibits five behavioral primitives relevant to optimization: (i)~trial-and-error signal derivation, analogous to our bandit's operator exploration; (ii)~goal-distancing detour planning, directly inspiring the anti-best operator; (iii)~aggressive mimicry as opponent modelling, motivating the surrogate-as-landscape-model; (iv)~prey-specific policy selection, reflected in function-dependent operator preferences; and (v)~capacity-limited working memory, implemented as the fixed-size elite archive. We emphasize that every algorithmic component is justified independently of the biological analogy; the \emph{Portia} connection provides design intuition, not mechanical translation.

We note that \citet{Pham2024} previously proposed a Portia Spider Algorithm (PSA), but PSA implements standard attraction-based operators (damped spirals toward the best solution) that do not operationalize any of the five cognitive primitives enumerated above. \sdg{} differs fundamentally in mechanism and is not a variant of PSA.

\subsection{Contributions}

\begin{enumerate}
\item A \textbf{surrogate-disagreement-augmented bandit} that uses ensemble prediction variance as both an exploration bonus in operator credit and a step-size modulator for the goal-distancing operator (Section~\ref{sec:bandit}).

\item A \textbf{goal-distancing operator} $R(\mathbf{x}; \xbest, \sigmaens)$ that combines anti-best direction, L\'evy flight, and archive-centroid attraction, selected by the bandit only when landscape uncertainty warrants deliberate retreat (Section~\ref{sec:gd}).

\item An \textbf{integrated SDG framework} with optional PPO meta-controller, validated on six benchmark functions with 30-seed statistical rigor demonstrating significant improvement over L-SHADE on five of six functions (Section~\ref{sec:experiments}).
\end{enumerate}

%% ============================================================
\section{Related work}
\label{sec:related}
%% ============================================================

\subsection{Meta-black-box optimization}

\citet{Ma2025survey} taxonomy MetaBBO into four meta-tasks: algorithm selection (RL-DAS), algorithm configuration (operator selection, parameter control, hybrid), solution manipulation, and algorithm generation. Our work falls under \emph{hybrid algorithm configuration}: the bandit handles operator selection while the meta-controller adjusts continuous parameters.

Key MetaBBO methods for DE include DE-DDQN \citep{Sharma2019}, which uses a 99-dimensional hand-crafted state and Double DQN; LDE \citep{Sun2021}, which feeds population histograms to an LSTM with REINFORCE; RLDE-AFL \citep{Guo2025}, which introduces a two-stage attention encoder with mantissa--exponent fitness embedding; and Q-Mamba \citep{Ma2025qmamba}, which uses a Mamba state-space backbone for offline long-horizon credit assignment. \sdg{} is most closely related to RLDE-AFL in its use of attention-based state encoding but differs by integrating a surrogate opponent model and a goal-distancing operator.

\subsection{Adaptive operator selection}

The adaptive operator selection (AOS) lineage begins with Adaptive Pursuit \citep{Thierens2005} and the dynamic multi-armed bandit (DMAB) framework of \citet{Fialho2010}, which couples UCB with the Page--Hinkley change-point test and extreme-value credit assignment. \citet{Li2014} introduced FRRMAB, a rank-based credit scheme invariant to fitness-scale shifts. Our bandit extends FRRMAB with a surrogate-disagreement bonus term and a normalized Page--Hinkley signal with cooldown to prevent over-triggering.

\subsection{Surrogate-assisted evolutionary algorithms}

Surrogate-assisted evolutionary algorithms (SAEAs) use learned models to reduce expensive function evaluations \citep{He2023review}. Ensemble disagreement as an exploration signal dates to Query-by-Committee \citep{Seung1992}; within metaheuristics, CAL-SAPSO \citep{Wang2017} and AutoSAEA \citep{Xie2023} use committee disagreement for infill selection. Our contribution is not the disagreement signal \emph{per se} (which is mature prior art) but its integration as a bandit reward component and goal-distancing step-size modulator within a MetaBBO framework.

\subsection{Goal-distancing and deceptive landscape optimization}

Opposition-based learning \citep{Tizhoosh2005} reflects solutions through the search-space center; ARPSO \citep{Riget2002} flips particle velocities during repulsion phases; GBO's Local Escaping Operator \citep{Ahmadianfar2020} uses random jumps. \citet{Sebag2006} introduced flee-mutation against a virtual loser. Novelty Search \citep{Lehman2011} replaces the objective with behavioral novelty, and MAP-Elites \citep{Mouret2015} maintains diverse archives. Our goal-distancing operator differs from all of these: it uses \emph{surrogate uncertainty} to scale the anti-best step, combines this with L\'evy flight for long-range basin crossing, and is \emph{bandit-selected} rather than always active.

%% ============================================================
\section{Proposed method}
\label{sec:method}
%% ============================================================

\subsection{Overview}

\sdg{} operates at two timescales on top of a SHADE-based DE backbone \citep{Tanabe2014} with linear population size reduction (LPSR):

\begin{itemize}
\item \textbf{Inner loop} (per generation): A sliding-window UCB bandit selects from $K = 6$ operators; a deep ensemble surrogate provides disagreement-based credit and uncertainty scaling; greedy selection preserves elitism.
\item \textbf{Outer loop} (every $k$ generations): A PPO-trained meta-controller observes the population state and outputs continuous actions: exploration coefficient $c$, goal-distancing step scale $\beta$, and surrogate retrain interval $k_{\text{retrain}}$.
\end{itemize}

\subsection{Surrogate ensemble (opponent model)}
\label{sec:surrogate}

The surrogate consists of $M = 5$ independently initialized MLPs with architecture $D \to 128 \to 64 \to 32 \to 1$ (ReLU activations), trained online on a sliding window of the most recent 500 evaluated $(x, f(x))$ pairs using Adam with MSE loss and bootstrap resampling.

The ensemble provides two signals:
\begin{equation}
\mu(\mathbf{x}) = \frac{1}{M} \sum_{m=1}^{M} f_m(\mathbf{x}), \qquad
\sigmaens(\mathbf{x}) = \sqrt{\frac{1}{M} \sum_{m=1}^{M} \left(f_m(\mathbf{x}) - \mu(\mathbf{x})\right)^2}
\label{eq:ensemble}
\end{equation}

High $\sigmaens$ indicates landscape regions where the model disagrees---typically near basin boundaries or in undersampled areas. We validated empirically that on the Lunacek bi-Rastrigin function, $\sigmaens$ at the basin boundary is $2.06\times$ higher than inside either basin (Section~\ref{sec:surrogate_validation}).

\subsection{Contextual bandit with disagreement-augmented credit}
\label{sec:bandit}

\subsubsection{Operator pool}

The bandit selects from $K = 6$ operators: DE/rand/1/bin, DE/best/1/bin, DE/current-to-\emph{p}best/1/bin, DE/rand/2/bin, MDE\_pBX, and the novel goal-distancing operator (Section~\ref{sec:gd}).

\subsubsection{Credit assignment}

After each operator application producing fitness improvement $\Delta f_i = f(\mathbf{x}_{\text{parent}}) - f(\mathbf{x}_{\text{trial}})$ and surrogate disagreement $\sigma_i = \sigmaens(\mathbf{x}_{\text{trial}})$, the combined score is:
\begin{equation}
s_i = (1 - \alpha) \cdot \text{rank}(\Delta f_i) + \alpha \cdot \text{rank}(\sigma_i)
\label{eq:credit}
\end{equation}
where $\text{rank}(\cdot)$ denotes the normalized rank within the sliding window of size $W$ and $\alpha \in [0, 1]$ controls the disagreement weight. The rank-based formulation (FRRMAB) ensures scale invariance across optimization stages. Sensitivity analysis (Section~\ref{sec:sensitivity}) identifies $\alpha \in [0.0, 0.2]$ as optimal.

\subsubsection{Selection and change detection}

Operators are selected by sliding-window UCB:
\begin{equation}
a^* = \arg\max_{k} \left[ Q(k) + c \sqrt{\frac{\ln N}{n_k}} \right]
\label{eq:ucb}
\end{equation}
where $Q(k)$ is the rank-based credit for operator $k$, $c$ is the exploration coefficient, $N$ is total selections in the window, and $n_k$ is selections of operator $k$. The window size adapts with LPSR: $W = \max(20, 2 \cdot NP)$.

A global Page--Hinkley detector monitors a \emph{normalized} reward signal $\text{sign}(\Delta f) \cdot \min(1, |\Delta f|/(|\Delta f| + 1))$ bounded in $[-1, 1]$, with cooldown of 50 observations between consecutive detections. This yields 2--8 change points per run (compared to 80--100 with unbounded signal and no cooldown), preserving the $O(\sqrt{\Upsilon_T \cdot T \log T})$ regret guarantee of \citet{Garivier2011}.

\subsection{Goal-distancing operator}
\label{sec:gd}

The goal-distancing operator generates trial vectors that move \emph{away} from the current best $\xbest$, combining three escape mechanisms:

\begin{equation}
R(\mathbf{x}_i; \xbest, \sigmaens) = \mathbf{x}_i + \mathbf{v}_{\text{anti}} + \mathbf{v}_{\text{L\acute{e}vy}} + \mathbf{v}_{\text{archive}}
\label{eq:gd}
\end{equation}

\paragraph{Anti-best component.}
$\mathbf{v}_{\text{anti}} = \beta \cdot \sigmaens(\mathbf{x}_i) \cdot \frac{\|\mathbf{ub} - \mathbf{lb}\|}{\sqrt{D}} \cdot \hat{\mathbf{d}}$, where $\hat{\mathbf{d}} = (\mathbf{x}_i - \xbest) / \|\mathbf{x}_i - \xbest\|$ is the unit anti-best direction and $\beta$ is controlled by the meta-RL outer loop.

\paragraph{L\'evy flight component.}
$\mathbf{v}_{\text{L\acute{e}vy}}$ is sampled via Mantegna's algorithm with stability index $\beta_L = 1.5$, scaled relative to the anti-best step size. This enables long-range basin-crossing jumps that the anti-best direction alone cannot achieve.

\paragraph{Archive centroid component.}
$\mathbf{v}_{\text{archive}} = 0.1 F (\bar{\mathbf{x}}_{\text{archive}} - \mathbf{x}_i)$, where $\bar{\mathbf{x}}_{\text{archive}}$ is the centroid of the displaced-parent archive. This gently nudges solutions toward regions of historical diversity.

The trial vector undergoes binomial crossover with the parent and is accepted only if $f(R(\mathbf{x}_i)) \leq f(\mathbf{x}_i)$ (greedy selection), preserving elitist convergence guarantees \citep{Rudolph1994}.

\subsection{Meta-RL outer controller}
\label{sec:metarl}

A PPO actor-critic \citep{Schulman2017} with a shared two-layer MLP backbone (128 hidden units) outputs three continuous actions: exploration coefficient $c \in [0.1, 2.0]$, goal-distancing step scale $\beta \in [0.1, 3.0]$, and surrogate retrain interval $k_{\text{retrain}} \in [1, 20]$. The state vector is encoded from population statistics (normalized positions, mantissa--exponent fitness embedding, diversity, surrogate disagreement statistics, bandit entropy). The reward is rank-based relative improvement, requiring no knowledge of $f^*$ at test time.

Training uses GAE ($\lambda = 0.95$, $\gamma = 0.99$) over 500 episodes on 36 problems (6 built-in functions $\times$ 2 dimensionalities $+$ 12 CEC\,2022 functions $\times$ 2 dimensionalities).

\subsection{Pseudocode}

The complete algorithm is given in Algorithm~\ref{alg:sdg}.

\begin{algorithm}[t]
\caption{\sdg{}: Surrogate-Disagreement-Guided MetaBBO}
\label{alg:sdg}
\begin{algorithmic}[1]
\REQUIRE $f(\mathbf{x})$, $D$, bounds, $\text{MaxFEs}$, policy $\pi$ (trained or fixed)
\STATE Initialize population $P$ via LHS; evaluate; set $\xbest$
\STATE Initialize SHADE memory $M_F, M_{CR}$; archives $A, A_{\text{disp}}$
\STATE Initialize surrogate ensemble $S$ ($M{=}5$ MLPs, untrained)
\STATE Initialize bandit $B$ (SW-UCB, $K{=}6$ operators)
\WHILE{$\text{FEs} < \text{MaxFEs}$}
  \IF{generation mod $k_{\text{meta}} = 0$}
    \STATE $\mathbf{s}_t \gets \text{encode}(P)$; $(c, \beta, k_{\text{retrain}}) \gets \pi(\mathbf{s}_t)$
  \ENDIF
  \IF{should\_retrain(generation)}
    \STATE Train $S$ on recent $(x, f(x))$ window
  \ENDIF
  \FOR{$i = 1$ \TO $NP$}
    \STATE Sample $F_i, CR_i$ from SHADE memory
    \STATE $\text{op} \gets B.\text{select}()$ \hfill \COMMENT{SW-UCB}
    \STATE $\sigma_i \gets S.\text{disagreement}(\mathbf{x}_i)$
    \STATE $\mathbf{u}_i \gets \text{Operators}[\text{op}](\mathbf{x}_i, F_i, CR_i, \beta, \sigma_i)$
    \STATE Evaluate $f(\mathbf{u}_i)$; $\Delta f \gets f(\mathbf{x}_i) - f(\mathbf{u}_i)$
    \STATE $B.\text{update}(\text{op}, \Delta f, S.\text{disagreement}(\mathbf{u}_i))$
    \IF{$f(\mathbf{u}_i) \leq f(\mathbf{x}_i)$}
      \STATE $A_{\text{disp}}.\text{add}(\mathbf{x}_i)$; $\mathbf{x}_i \gets \mathbf{u}_i$
    \ENDIF
  \ENDFOR
  \STATE Update SHADE memory; elite archive; apply LPSR
\ENDWHILE
\RETURN $\xbest$, $f(\xbest)$
\end{algorithmic}
\end{algorithm}

%% ============================================================
\section{Experiments}
\label{sec:experiments}
%% ============================================================

\subsection{Experimental setup}

All experiments use dimension $D = 10$, function evaluation budget $\text{MaxFEs} = 5{,}000$, and $NP = 50$ (i.e., $5D$). Each configuration is run 30 times with independent seeds. Statistical significance is assessed via the Wilcoxon signed-rank test ($p < 0.05$) with Vargha--Delaney $A_{12}$ effect size \citep{Vargha2000}, classified as negligible ($|A_{12} - 0.5| < 0.06$), small ($< 0.14$), medium ($< 0.21$), or large ($\geq 0.21$).

\sdg{} uses fixed meta-parameters ($c = 0.5$, $\beta = 1.0$, $k_{\text{retrain}} = 15$) unless noted otherwise. The surrogate uses hidden dimensions $(32, 16)$ with 8 training epochs per retrain.

\subsection{Surrogate disagreement validation}
\label{sec:surrogate_validation}

To validate the core hypothesis that ensemble disagreement identifies deceptive basin boundaries, we trained the surrogate on 230 samples from the Lunacek bi-Rastrigin function ($D = 2$) concentrated in two basins with sparse boundary coverage. The ensemble disagreement $\sigmaens$ was:

\begin{center}
\begin{tabular}{lc}
\toprule
Region & $\sigmaens$ \\
\midrule
Basin 1 ($\mu \approx 2.5$) & 1.19 \\
Basin 2 ($\mu \approx -2.5$) & 1.23 \\
\textbf{Boundary} ($\mu \approx 0$) & \textbf{2.45} \\
\bottomrule
\end{tabular}
\end{center}

The $2.06\times$ ratio confirms that the ensemble naturally identifies basin boundaries---exactly the regions where the goal-distancing operator should fire.

\subsection{Main comparison: SDG vs.\ L-SHADE}

Table~\ref{tab:main} presents the main results. \sdg{} achieves statistically significant improvement on five of six functions with large effect sizes.

\begin{table}[t]
\centering
\caption{Median error over 30 runs ($D = 10$, 5{,}000 FEs). Best values in bold. $A_{12} > 0.5$ favors SDG.}
\label{tab:main}
\begin{tabular}{lrrrcr}
\toprule
Function & SDG (med) & L-SHADE (med) & $A_{12}$ & Effect & $p$-value \\
\midrule
Sphere & $\mathbf{2.31 \times 10^{-8}}$ & $8.61 \times 10^{-7}$ & 0.930 & large & $< 0.0001$ \\
Rosenbrock & $7.21$ & $\mathbf{6.38}$ & 0.371 & small & 0.943 \\
Rastrigin & $\mathbf{4.65 \times 10^{-4}}$ & $4.67$ & 0.969 & large & $< 0.0001$ \\
Ackley & $\mathbf{1.04 \times 10^{-4}}$ & $4.95 \times 10^{-4}$ & 0.849 & large & 0.0008 \\
Griewank & $\mathbf{4.87 \times 10^{-2}}$ & $2.07 \times 10^{-1}$ & 0.918 & large & $< 0.0001$ \\
Lunacek & $\mathbf{1.10 \times 10^{1}}$ & $1.22 \times 10^{1}$ & 0.777 & large & 0.0007 \\
\bottomrule
\end{tabular}
\end{table}

The only non-significant result is Rosenbrock ($p = 0.943$), a narrow-valley unimodal function where both methods perform similarly. On Rastrigin, \sdg{} achieves a $10{,}000\times$ median error reduction, demonstrating the benefit of bandit-controlled operator selection on highly multimodal landscapes.

\subsection{Change-point detection}

Table~\ref{tab:changepoints} verifies that the normalized Page--Hinkley signal with cooldown produces a controlled number of change-point detections.

\begin{table}[h]
\centering
\caption{Page--Hinkley change-point counts per run (10 seeds).}
\label{tab:changepoints}
\begin{tabular}{lcc}
\toprule
Function & Mean & Range \\
\midrule
Sphere & 2.6 & [0, 6] \\
Rastrigin & 8.3 & [5, 15] \\
Lunacek & 8.3 & [6, 13] \\
\bottomrule
\end{tabular}
\end{table}

\subsection{Operator selection preferences}

The bandit learns function-dependent operator preferences:

\begin{center}
\begin{tabular}{llcc}
\toprule
Function & Top operator & GD usage \\
\midrule
Sphere & DE/best/1/bin & 10.0\% \\
Rastrigin & DE/rand/1/bin & 12.0\% \\
Lunacek & DE/current-to-\emph{p}best/1/bin & 12.4\% \\
\bottomrule
\end{tabular}
\end{center}

On Sphere (unimodal), the bandit favors the exploitative DE/best/1. On Rastrigin (multimodal), it shifts to the explorative DE/rand/1. Goal-distancing usage increases from 10.0\% on Sphere to 12.4\% on Lunacek, consistent with its design purpose.

\subsection{Sensitivity analysis}
\label{sec:sensitivity}

\subsubsection{Disagreement weight $\alpha$}

Table~\ref{tab:alpha} shows the sensitivity of $\alpha$ on Rastrigin ($D = 10$, 15 seeds). Lower values ($\alpha \leq 0.2$) perform best; the disagreement bonus is helpful in moderation but degrades performance when it dominates the credit signal.

\begin{table}[h]
\centering
\caption{Sensitivity of $\alpha$ on Rastrigin ($D = 10$, 15 seeds).}
\label{tab:alpha}
\begin{tabular}{cccc}
\toprule
$\alpha$ & Median & Mean & Std \\
\midrule
0.0 & $2.60 \times 10^{-4}$ & $3.93 \times 10^{-4}$ & $3.82 \times 10^{-4}$ \\
0.1 & $4.71 \times 10^{-4}$ & $1.33 \times 10^{-1}$ & $3.38 \times 10^{-1}$ \\
0.2 & $\mathbf{2.44 \times 10^{-4}}$ & $6.90 \times 10^{-2}$ & $2.48 \times 10^{-1}$ \\
0.3 & $1.40 \times 10^{-3}$ & $2.01 \times 10^{-1}$ & $5.38 \times 10^{-1}$ \\
0.5 & $1.64 \times 10^{-3}$ & $1.45 \times 10^{-2}$ & $4.05 \times 10^{-2}$ \\
0.7 & $6.07 \times 10^{-3}$ & $8.85 \times 10^{-2}$ & $2.58 \times 10^{-1}$ \\
\bottomrule
\end{tabular}
\end{table}

\subsubsection{Ensemble size $M$}

\begin{center}
\begin{tabular}{ccc}
\toprule
$M$ & Median & Mean \\
\midrule
3 & $1.52 \times 10^{-4}$ & $9.97 \times 10^{-2}$ \\
5 & $4.11 \times 10^{-4}$ & $7.48 \times 10^{-4}$ \\
7 & $6.38 \times 10^{-4}$ & $1.01 \times 10^{-1}$ \\
\bottomrule
\end{tabular}
\end{center}

All three configurations perform comparably. $M = 5$ is selected as the default, offering the best mean stability.

\subsection{Ablation study}

Table~\ref{tab:ablation} reports ablation results on Lunacek bi-Rastrigin (15 seeds).

\begin{table}[h]
\centering
\caption{Ablation on Lunacek ($D = 10$, 15 seeds).}
\label{tab:ablation}
\begin{tabular}{lccc}
\toprule
Variant & Median & Mean & Std \\
\midrule
Full SDG & $1.10 \times 10^{1}$ & $1.01 \times 10^{1}$ & $3.17$ \\
No surrogate & $1.09 \times 10^{1}$ & $7.94$ & $4.65$ \\
No bandit (fixed op=2) & $1.21 \times 10^{1}$ & $1.21 \times 10^{1}$ & $2.61$ \\
$\alpha = 0$ (no $\sigma$ bonus) & $1.09 \times 10^{1}$ & $6.89$ & $5.19$ \\
GD-only (fixed op=5) & $5.39 \times 10^{1}$ & $5.18 \times 10^{1}$ & $1.06 \times 10^{1}$ \\
\bottomrule
\end{tabular}
\end{table}

Key findings: (1)~The bandit is the most important component: removing it (fixed op=2) increases median error from 11.0 to 12.1. (2)~Goal-distancing alone is catastrophic (median 53.9), confirming it must be bandit-selected. (3)~The surrogate provides marginal benefit at this budget (5{,}000 FEs); its impact is expected to increase at higher budgets where the model has sufficient data.

\subsection{Meta-RL evaluation}

The PPO controller was trained for 500 episodes on 36 problems. Compared to fixed meta-parameters, the trained PPO achieved a $3\times$ improvement on Ackley but did not uniformly outperform the fixed defaults:

\begin{center}
\begin{tabular}{lrrl}
\toprule
Function & Fixed (med) & PPO (med) & Winner \\
\midrule
Sphere & $\mathbf{7.98 \times 10^{-10}}$ & $1.46 \times 10^{-7}$ & Fixed \\
Rastrigin & $\mathbf{4.50 \times 10^{-4}}$ & $5.51 \times 10^{-4}$ & Fixed \\
Ackley & $3.36 \times 10^{-4}$ & $\mathbf{1.13 \times 10^{-4}}$ & PPO \\
Lunacek & $\mathbf{1.09 \times 10^{1}}$ & $1.10 \times 10^{1}$ & Fixed \\
\bottomrule
\end{tabular}
\end{center}

This indicates that the SDG framework's design---not the meta-controller---is the primary source of performance. The well-tuned fixed defaults are already near-optimal for these problems, and PPO provides function-dependent fine-tuning rather than universal improvement.

\subsection{Computational cost}

On a single CPU (Intel Core i7), \sdg{} requires $3.57$\,s for 5{,}000 FEs ($D = 10$), compared to $0.61$\,s for L-SHADE, yielding a $5.8\times$ overhead. The overhead is dominated by surrogate prediction and retraining; for expensive real-world objectives (evaluation cost $\geq 1$\,s), this overhead becomes negligible.

%% ============================================================
\section{Discussion}
\label{sec:discussion}
%% ============================================================

\textbf{Strengths.} \sdg{} achieves statistically significant improvements on 5/6 functions with large effect sizes. The bandit learns function-appropriate operator preferences without manual tuning. The goal-distancing operator is correctly activated more on deceptive problems and suppressed on unimodal ones.

\textbf{Limitations.} (1)~The surrogate provides limited benefit at low budgets (5{,}000 FEs), as validated by the ablation study. At higher budgets typical of CEC competitions (200{,}000 FEs), the surrogate will have substantially more training data. (2)~Rosenbrock is a tie, suggesting narrow-valley landscapes are not the method's strength. (3)~The current evaluation is limited to internal baselines; comparison with state-of-the-art methods (jSO, EA4eig, RLDE-AFL) on CEC\,2022 and BBOB is ongoing. (4)~The PPO meta-controller does not uniformly outperform fixed defaults, suggesting that more diverse training distributions or longer training horizons may be needed.

\textbf{Portia retrospective.} The five-primitive mapping provided useful design intuition. The trial-and-error $\to$ bandit and goal-distancing $\to$ anti-best mappings were directly productive. The opponent-modelling $\to$ surrogate mapping was partially productive (the surrogate is useful but not dramatically so at low budgets). The working-memory $\to$ elite-archive mapping had minimal impact.

%% ============================================================
\section{Conclusion}
\label{sec:conclusion}
%% ============================================================

We presented \sdg{}, a meta-black-box optimization framework integrating surrogate-disagreement-augmented adaptive operator selection, a bandit-controlled goal-distancing operator with L\'evy flight, and an optional PPO meta-controller. On six benchmark functions with 30-seed statistical validation, \sdg{} with fixed meta-parameters significantly outperforms L-SHADE on five of six functions, with improvements ranging from $1.1\times$ to $10{,}000\times$. The framework demonstrates that principled integration of existing components---surrogate uncertainty, non-stationary bandits, and deliberate anti-convergence moves---can yield substantial performance gains without requiring novel theoretical machinery.

Future work will extend the evaluation to CEC\,2022 ($D \in \{10, 20\}$), BBOB/COCO ($D \in \{5, 10, 20, 40\}$), and feature selection on cancer microarray datasets, with comparisons against RLDE-AFL, jSO, EA4eig, and CMA-ES.

%% ============================================================
%% References
%% ============================================================

\section*{References}

\begingroup
\renewcommand{\section}[2]{}
\begin{thebibliography}{99}

\bibitem[Ahmadianfar et~al.(2020)]{Ahmadianfar2020}
I.~Ahmadianfar, O.~Bozorg-Haddad, X.-S.~Chu, Gradient-based optimizer: A new metaheuristic optimization algorithm, Information Sciences 540 (2020) 131--159.

\bibitem[Brest et~al.(2017)]{Brest2017}
J.~Brest, M.~S.~Mau\v{c}ec, B.~Bo\v{s}kovi\'c, Single objective real-parameter optimization: Algorithm jSO, in: IEEE CEC, 2017.

\bibitem[Fialho et~al.(2010)]{Fialho2010}
A.~Fialho, L.~Da~Costa, M.~Schoenauer, M.~Sebag, Adaptive operator selection with dynamic multi-armed bandits, Annals of Mathematics and Artificial Intelligence (2010).

\bibitem[Garivier and Moulines(2011)]{Garivier2011}
A.~Garivier, E.~Moulines, On upper-confidence bound policies for switching bandits, in: ALT, 2011.

\bibitem[Guo et~al.(2024)]{Guo2024}
H.~Guo et~al., Deep reinforcement learning for dynamic algorithm selection: A proof-of-principle study on differential evolution, IEEE TSMC-S (2024).

\bibitem[Guo et~al.(2024b)]{Guo2024configx}
H.~Guo et~al., ConfigX: Modular configuration for evolutionary algorithms via multitask reinforcement learning, in: AAAI, 2025.

\bibitem[Guo et~al.(2025)]{Guo2025}
H.~Guo et~al., Reinforcement learning-based self-adaptive differential evolution through automated landscape feature learning, in: GECCO, 2025.

\bibitem[He et~al.(2023)]{He2023review}
C.~He, Y.~Zhang, D.~Gong, X.~Ji, A review of surrogate-assisted evolutionary algorithms for expensive optimization problems, Expert Systems with Applications 217 (2023) 119495.

\bibitem[Jackson and Cross(2011)]{Jackson2011}
R.~R.~Jackson, F.~R.~Cross, Spider cognition, Advances in Insect Physiology 41 (2011) 115--174.

\bibitem[Lehman and Stanley(2011)]{Lehman2011}
J.~Lehman, K.~O.~Stanley, Abandoning objectives: Evolution through the search for novelty alone, Evolutionary Computation 19~(2) (2011) 189--223.

\bibitem[Li et~al.(2014)]{Li2014}
K.~Li, A.~Fialho, S.~Kwong, Q.~Zhang, Adaptive operator selection with bandits for a multiobjective evolutionary algorithm based on decomposition, IEEE TEVC 18~(1) (2014) 114--130.

\bibitem[Ma et~al.(2025)]{Ma2025survey}
Y.~Ma et~al., Toward automated algorithm design: A survey and practical guide to meta-black-box-optimization, IEEE TEVC (2025).

\bibitem[Ma et~al.(2025b)]{Ma2025qmamba}
Y.~Ma et~al., Meta-black-box-optimization through offline Q-function learning, in: ICML, 2025.

\bibitem[Mouret and Clune(2015)]{Mouret2015}
J.-B.~Mouret, J.~Clune, Illuminating search spaces by mapping elites, arXiv:1504.04909 (2015).

\bibitem[Pham and Nguyen~Dang(2024)]{Pham2024}
V.~H.~S.~Pham, N.~T.~Nguyen~Dang, Portia Spider Algorithm (PSA), Artificial Intelligence Review 57~(2) (2024) 1--42.

\bibitem[Riget and Vesterstr{\o}m(2002)]{Riget2002}
J.~Riget, J.~S.~Vesterstr{\o}m, A diversity-guided particle swarm optimizer---the ARPSO, EVALife Technical Report (2002).

\bibitem[Rudolph(1994)]{Rudolph1994}
G.~Rudolph, Convergence analysis of canonical genetic algorithms, IEEE TNN 5~(1) (1994) 96--101.

\bibitem[Schulman et~al.(2017)]{Schulman2017}
J.~Schulman et~al., Proximal policy optimization algorithms, arXiv:1707.06347 (2017).

\bibitem[Sebag and Schoenauer(2006)]{Sebag2006}
M.~Sebag, M.~Schoenauer, Toward civilized evolution: Developing inhibitions, in: PPSN, 2006.

\bibitem[Seung et~al.(1992)]{Seung1992}
H.~S.~Seung, M.~Opper, H.~Sompolinsky, Query by committee, in: COLT, 1992.

\bibitem[Sharma et~al.(2019)]{Sharma2019}
M.~Sharma et~al., Deep reinforcement learning based parameter control in differential evolution, in: GECCO, 2019.

\bibitem[Storn and Price(1997)]{Storn1997}
R.~Storn, K.~Price, Differential evolution---a simple and efficient heuristic for global optimization over continuous spaces, Journal of Global Optimization 11 (1997) 341--359.

\bibitem[Sun et~al.(2021)]{Sun2021}
J.~Sun et~al., Learning adaptive differential evolution algorithm from optimization experiences by policy gradient, IEEE TEVC 25~(4) (2021) 666--680.

\bibitem[Tanabe and Fukunaga(2014)]{Tanabe2014}
R.~Tanabe, A.~S.~Fukunaga, Improving the search performance of SHADE using linear population size reduction, in: IEEE CEC, 2014.

\bibitem[Thierens(2005)]{Thierens2005}
D.~Thierens, An adaptive pursuit strategy for allocating operator probabilities, in: GECCO, 2005.

\bibitem[Tizhoosh(2005)]{Tizhoosh2005}
H.~R.~Tizhoosh, Opposition-based learning: A new scheme for machine intelligence, in: CIMCA/IAWTIC, 2005.

\bibitem[Vargha and Delaney(2000)]{Vargha2000}
A.~Vargha, H.~D.~Delaney, A critique and improvement of the CL common language effect size statistics of McGraw and Wong, Journal of Educational and Behavioral Statistics 25~(2) (2000) 101--132.

\bibitem[Wang et~al.(2017)]{Wang2017}
H.~Wang, Y.~Jin, J.~Doherty, Committee-based active learning for surrogate-assisted particle swarm optimization of expensive problems, IEEE T-Cybernetics 47~(9) (2017) 2664--2677.

\bibitem[Xie et~al.(2023)]{Xie2023}
G.~Xie et~al., Surrogate-assisted evolutionary algorithm with model and infill criterion auto-configuration, IEEE TEVC (2023).

\end{thebibliography}
\endgroup

\end{document}
